\chapter{Reliable cross-system integration}
\label{ch:enterprise-outbox}
\chapterbrief{Record durable intent and safely replay delivery.}
\par\smallskip\noindent\textbf{Abbreviations used in this chapter:} DBI --- Database Interface; HTTP --- Hypertext Transfer Protocol; SQL --- Structured Query Language; ID --- identifier.\par\smallskip

\index{Transactional outbox}
\index{Inbox}
\index{Idempotency}
\index{Replay}

A local SQL transaction cannot roll back an HTTP request or a message
already accepted by another system. Holding a local transaction open while
calling the external system does not solve this problem; it also increases
lock duration. Instead, separate durable intent from delivery.

\strategy{app-outbox}{Transactional outbox}{Write the shipping instruction on the same receipt as the sale.}{a committed change must trigger a remote effect.}{the effect can be committed in the same local transaction.}{DBI, unique keys, durable queues}

\section{Transactional outbox and inbox}
\index{Transactional outbox}\index{Idempotency}
\begin{figure}[!htbp]\centering
\begin{tikzpicture}[node distance=10mm]
\node[okbox,text width=5cm] (s) {Source database\\business row + outbox event\\one transaction};
\node[box,right=of s,text width=3cm] (r) {Relay\\claim and send};
\node[okbox,below=of r,text width=3cm] (d) {Destination\\unique event ID};
\node[box,left=of d,text width=5cm] (a) {Source acknowledgement\\conditional on claim};
\draw[arr](s)--(r); \draw[arr](r)--(d) node[midway,lbl,right]{stable ID};
\draw[arr](d)--(a); \draw[arr](a)--(s);
\node[badbox,below=of a,text width=5cm] (f) {Crash after remote commit?\\Retry with the same ID};
\draw[badarr](f.east)-|(d.south);
\end{tikzpicture}
\caption{The outbox preserves intent; destination deduplication resolves ambiguous delivery.}
\end{figure}
Insert an outbox event alongside the business row in the same transaction
[71]. A relay claims committed events, delivers them, and records completion
after acknowledgement. If the response is lost, the relay cannot know
whether the remote operation happened. Reuse the event ID on retries.
The receiving application should commit its deduplication record (inbox)
and effect in a single transaction. A unique inbox row inserted before a
nontransactional external effect does not, by itself, guarantee exactly
one effect.

\section{Failure windows and recovery}
\begin{longtable}{L{.34\linewidth}L{.58\linewidth}}
\toprule Failure point & Recovery \\ \midrule\endhead
Before source commit & Neither business state nor intent is committed; replay input. \\
After source commit & Outbox remains pending; relay later. \\
After destination commit, before acknowledgement & Retry same event ID; destination returns prior result or deduplicates. \\
After acknowledgement & No delivery needed; retain audit and dedupe records for the replay horizon. \\
\bottomrule
\end{longtable}
External APIs vary: some support idempotency keys, others offer lookup by a
client reference, and some offer neither. For the last case, ambiguous
outcomes may require reconciliation or human review. Compensation is a new
business operation, not a rollback of history. Preserve ordering by
aggregate when required, and record failures with correlation identifiers
rather than credentials or full sensitive payloads.

\section{Laboratory: two independent databases}
The following laboratory uses separate local SQLite databases, one for
business state and outbox and one for the receiving effect. It injects a
failure after the destination commits but before the source acknowledges,
then replays the same event. It also verifies source rollback and rejects
reuse of an input identity with a different value. The relay is intentionally
single-worker; multiple relays need the claim protocol of
Chapter~\ref{ch:enterprise-pipelines}.
\CodeFile{perl}{examples/34-replay-outbox.pl}{Example 34: rollback, replay and destination deduplication}

